\begin{center}
\begin{tikzpicture}[>=stealth, xscale=0.44, yscale=3]
\fill[blue!25, domain=6:12, samples=40] (6,0) -- plot (\x, {0.36675*\x*\x*exp(-((\x/3.948)^2))}) -- (12,0) -- cycle;

 \draw[gray!20, xstep=1, ystep=0.1] (0,0) grid (25,2.2);
 \draw[gray!50, xstep=5, ystep=0.5] (0,0) grid (25,2.2);
 \draw[->] (0,0) -- (26.5,0) node[below left=12pt] {$v$ in \si{\meter\per\second}};
 \draw[->] (0,0) -- (0,2.45) node[above right] {$f(v)$ in $10^{-3}\,\mathrm{s/m}$};
 \foreach \x/\l in {0/0,5/500,10/1000,15/1500,20/2000,25/2500} \node[below] at (\x,0) {\l};
 \foreach \y in {0.5,1,1.5,2} \node[left] at (0,\y) {\y};
 \draw[very thick, blue, domain=0:12, samples=80] plot (\x, {0.36675*\x*\x*exp(-((\x/3.948)^2))});
 \draw[very thick, red, domain=0:22, samples=80] plot (\x, {0.060249*\x*\x*exp(-((\x/7.2085)^2))});
 \draw[very thick, black!70!green, domain=0:25, samples=80] plot (\x, {0.021303*\x*\x*exp(-((\x/10.194)^2))});
 \node[blue, above right] at (4.2,2.05) {\TAO};
 \node[red, above right] at (7.6,1.12) {\TBO};
 \node[black!70!green, above right] at (13.5,0.72) {\TCO};

 \draw[|<->|, thick, blue] (1.90,1.05) -- (6.46,1.05);
 \draw[|<->|, thick, red] (3.47,0.575) -- (11.80,0.575);
 \draw[|<->|, thick, black!70!green] (4.91,0.407) -- (16.68,0.407);
 \draw[dashed, blue] (3.948,0) -- (3.948,2.10);
\end{tikzpicture}
\end{center}

\begin{abcliste}
\abc Die wahrscheinlichste Geschwindigkeit liegt beim Maximum der Kurve. Man liest sie beim Maximum der blauen Kurve ab. Genauer: Nach FoTaBe S.~171 ist $v_\text{w} = \sqrt{2kT/m} = \sqrt{2RT/M}$, mit $M = \MOzF = 2\cdot\MOaO$:
\begin{align}
 v_\text{w} &= \vwAF = \sqrt{\frac{2\cdot\ncR\cdot\TA}{\MOz}} = \vwA \approx \result{\vwAP{0}{2}}
\end{align}

\abc Man zeichnet bei halber Höhe des Maximums eine waagrechte Strecke zwischen den beiden Flanken der Kurve und liest ihre Länge ab:
\begin{align}
 \Delta v_{300} &\approx \result{\bArP{0}{2}}, & \Delta v_{1000} &\approx \result{\bBrP{0}{2}}, & \Delta v_{2000} &\approx \result{\bCrP{0}{3}}
\end{align}
Die Breite wächst wie $v_\text{w}$ proportional zu $\sqrt{T}$: Bei höherer Temperatur ist die Verteilung breiter und flacher.

\abc Der gesuchte Bruchteil ist die Fläche unter der Kurve rechts von \vgrO\ (blau), verglichen mit der ganzen Fläche unter der Kurve, die allen Molekülen entspricht. Ein kleines Gitterfeld ($\dvGO\cdot\dfGO$) entspricht $\aGP{0}{1}$ der Moleküle. Das Auszählen der Felder ergibt
\begin{align}
 \frac{\Delta N}{N} &\approx \result{\anrrP{0}{2}}
\end{align}
\end{abcliste}
